\section{The Translation}\label{sec:translation}

In this section, we define two languagues, a base language \(\LB\) that
corresponds a simple functional language~(\cref{subsec:base-lang}) and a logic
language $L$ that corresponds to a functional logic
language~(\cref{subsec:logic-lang}). We then develop a \emph{clairvoyance
  translation} strategy that translates a program in \(\LB\) to a program in
\(\LL\) that models the clairvoyance
semantics~\cite{clairvoyant,forking-paths}~(\cref{subsec:clairvoyance-translation}).
We then develop a \emph{demand translation}, which models the demand
semantics~\cite{bjerner1989,demand-semantics}, as a thin wrapper of the
clairvoyance translation~(\cref{subsec:demand-translation}). Finally, we show
how such a translation strategy works by applying it to the example of
\texttt{take} and \texttt{append} discussed in \cref{sec:motivating-example}.

We prove that this translation strategy is correct with respect to the
clairvoyance semantics and the demand semantics in Agda, but we defer details of
the correctness theorem, its proof, and the Adga development to
\cref{sec:equiv}.

\subsection{The Base Language \(\LB\)}\label{subsec:base-lang}

We show the syntax of \(\LB\) in \cref{fig:syntax-base}. The base language is a
simple typed calculus with booleans, lists, and lazy thunks. Its terms include
variables, let-bindings, conditionals, and recursions based on $\tfoldrZ$. In
addition, it contains the explicit \( \tlazyZ \) and \( \tforceZ \) for wrapping
a computation in a lazy thunk and forcing a computation in a thunk,
respectively. Conceptually, a thunk only evaluates the first time it is
\( \tforceZ \)d, then it remains that value.\nc{This would be true in an
  operational semantics, but it is not really an accurate description of any of
  the semantics in this paper.} Finally, we have an explicit \( \ttickZ \) term
to indicate incurring one unit of computation cost. These explicit operators
decouple our analysis from any particular evaluation model.

The language does not have functions or general recursions. This is because
(1)~the demand semantics does not support those
features~\cite{bjerner1989,demand-semantics}, and (2)~most functional logic
programming languages, with the exception of the Verse calculus~\cite{verse},
only support higher-order functions via defunctionalization~\cite[Section
3.3]{curry-lang}.

\begin{figure}[t]
\begin{center}
\begin{bnf}
  \( A \) ::= \TBool{} // \TList{A} // \TThunk{A} ;;

  $t$ ::= $x$ // \tlet{x}{t}{t} // \ttrue{} // \tfalse{} // \tif{t}{t}{t} |
  \tnil{} // \tcons{t}{t} // \tfoldr{x}{x}{t}{t}{t} // \tlazy{t} // \tforce{t} //
  \ttick{t}
\end{bnf}
\end{center}
\caption{The syntax of the base language \(\LB\).}\label{fig:syntax-base}
\end{figure}

\begin{figure}[t]
  \begin{center}
    \begin{mathpar}
      \inferrule{\Gamma \vdash^\LB t_1 : A \\ \Gamma \vdash^\LB t_2 :
        \TThunk{(\TList{A})}}{\Gamma \vdash^\LB \tcons{t_1}{t_2} : \TList{A}}{\textsc{cons}} \\
      \inferrule{\Gamma, x_1 : A, x_2 : \TThunk{B} \vdash^\LB t_1 : B \\ \Gamma
        \vdash^\LB t_2 : B \\ \Gamma \vdash^\LB t_3 : \TList{A}}{\Gamma \vdash^\LB \tfoldr{x_1}{x_2}{t_1}{t_2}{t_3} : B}{\textsc{foldr}} \\
      \inferrule{\Gamma \vdash^\LB t : A}{\Gamma \vdash^\LB \tlazy{t} :
        \TThunk{A}}{\textsc{lazy}} \qquad \inferrule{\Gamma \vdash^\LB t :
        \TThunk{A}}{\Gamma \vdash^\LB \tforce{t} : A}{\textsc{force}} \qquad
      \inferrule{\Gamma \vdash^\LB t : A}{\Gamma \vdash^\LB \ttick{t} :
        A}{\textsc{tick}}
    \end{mathpar}
  \end{center}
  \caption{Typing rules in $B$ that are related to lazy thunks and cost.}\label{fig:typing-base}
\end{figure}

Most of the typing rules for \(\LB\) are standard, so we omit them. We show the
typing rules for lazy thunks and computation costs in \cref{fig:typing-base}.
Notably, our lists are lazy: the cons of lists $\tcons{t_1}{t_2}$ takes a term
$t_1$ of type $A$ and a term $t_2$ of type $\TThunk{(\TList{A})}$, so the tail
of a list must always be wrapped in a thunk. Similarly, the ``function''
argument to $\tfoldrZ$ always assumes that its second paramter is wrapped in a
thunk.

\paragraph{Evaluation semantics.}
We can give language \(\LB\) a denotational semantics $\EE{t}{\gamma}$ for any
term $t$ under an evaluation context $\gamma$ by simply ignoring all
lazy-evaluation operators or computation cost:
$$\EE{\tlazy{t}}{\gamma} = \EE{\tforce{t}}{\gamma} = \EE{\ttick{t}}{\gamma} = \EE{t}{\gamma}$$

\subsection{The Logic Language $L$}\label{subsec:logic-lang}

We show the syntax of \(\LL\) in \cref{fig:syntax-logic}. Relative to \(\LB\),
\(\LL\) has some new types: a unit type \(\TUnit\), product types
\(\TProd{A}{B}\), and natural numbers \(\TNat\). We need the natural number type
to model computation costs, and we need product types to return computation
results along with costs. Since we are going to model demand semantics, we also
need products to model evaluation contexts, and we need the unit type to model
\emph{empty} contexts.

At the term level, \(\LL\) has the usual introduction and elimination forms for
the new types. We have no need for arbitrary functions on natural numbers, so we
include only literals and addition. We also have constructs for logic
programming: logic variables (\(\tfreeZ\)), nondeterministic choice
(\(\tchooseZ\)), and failure (\(\tfail\)). Note that logic variables must be
annotated with their types in order to ensure deterministic typing. Finally, we
have a new operator \(\tleZ\) that indicates whether its left-hand operand lies
below its right-hand operand in the order of approximations.

\begin{figure}[t]
\begin{center}
\begin{bnf}
  \( A \) ::=
  | \TUnit
  // \TBool
  // \TProd{A}{A}
  // \TThunk{A}
  // \TNat
  // \TList{A}
  ;;

  $t$ ::=
  | $x$
  // \tlet{x}{t}{t}
  // \tunit{}
  // \ttrue{}
  // \tfalse{}
  // \teq{t}{t}
  // \tle{t}{t}
  // \tif{t}{t}{t}
  | \tpair{t}{t}
  // \tfst{t}
  // \tsnd{t}
  | \tthunk{t}
  // \tundefined{}
  // \tcase{t}{x}{t}{t}
  | $n$
  // \tadd{t}{t}
  // \tnil{}
  // \tcons{t}{t}
  // \tfoldr{x}{x}{t}{t}{t}
  // \tchoose{t}{t}
  // \tfree{A}
  // \tfail{}
\end{bnf}
\end{center}
\caption{The syntax of the logic language \(\LL\).}\label{fig:syntax-logic}
\end{figure}

We define a simple functional logic language \(\LL\) instead of using an
existing formalization of one specific functional logic
language~(\eg~Curry~\cite{curry-lang,curry-monadic-semantics}, the Verse
calculus~\cite{verse}) to keep our approach generally applicable.\ly{Is there a
  particular formalization from a prior work that we are similar to?}

\paragraph{Typing.}

\ly{Why is there only one rule? Are other rules mostly straightforward? Or you
  haven't got the time to write them down?}

\begin{center}
  \begin{mathpar}
    \inferrule{\TB{\Gamma, x_1 : A, x_2 : \TThunk{B}}{t_1}{B} \\
      \TB{\Gamma}{t_2}{B} \\
      \TB{\Gamma}{t_3}{\TList{A}}}
    {\TB{\Gamma}{\tfoldr{x_1}{x_2}{t_1}{t_2}{t_3}}{B}}
  \end{mathpar}
\end{center}

\paragraph{Semantics.}

We show the semantic rules of $L$ that are related to the \TThunkZ{} datatype, the
choice operator~(\texttt{?}), and free variables in
\cref{fig:logic-semantics}.\ly{Did we refer to any papers when we define this
  semantics? If so, cite them.}

Note that $L$ is not total due to \( \tfail \).

\begin{figure}[t]
\begin{center}
  \begin{mathpar}
    \inferrule{\EL{\gamma}{t}{v}}{\EL{\gamma}{\tthunk{t}}{\sthunk{v}}}{\rulename{L-Result}}
    \qquad
    \inferrule{\\}{\EL{\gamma}{\tundefined}{\sundefined}}{\rulename{L-Undefined}} \\
    \inferrule{\EL{\gamma}{t_1}{\sthunk{v_1}} \\ \EL{\gamma, x \mapsto v_1}{t_2}{v_2}}{\EL{\gamma}{\tcase{t_1}{x}{t_2}{t_3}}{v_2}}{\rulename{L-Case-Result}} \\
    \inferrule{\EL{\gamma}{t_1}{\sundefined} \\ \EL{\gamma}{t_3}{v}}{\EL{\gamma}{\tcase{t_1}{x}{t_2}{t_3}}{v}}{\rulename{L-Case-Undefined}}\\
    \inferrule{v : A}{\EL{\gamma}{\tfree{A}}{v}}{\rulename{L-Free}} \qquad
    \inferrule{\EL{\gamma}{t_1}{v}}{\EL{\gamma}{\tchoose{t_1}{t_2}}{v}}{\rulename{L-Choice-L}}
    \qquad
    \inferrule{\EL{\gamma}{t_2}{v}}{\EL{\gamma}{\tchoose{t_1}{t_2}}{v}}{\rulename{L-Choice-R}} \\
    \inferrule{\EL{\gamma}{t_3}{v_1} \\
      \ELF{\gamma}{x_1}{x_2}{t_1}{t_2}{v_1}{v_2}}{\EL{\gamma}{\tfoldr{x_1}{x_2}{t_1}{t_2}{t_3}}{v_2}}{\rulename{L-Foldr}} \\
    \inferrule{\EL{\gamma}{t_2}{v}}{\ELF{\gamma}{x_1}{x_2}{t_1}{t_2}{\snil}{v}}{\rulename{L-Foldr-Nil}} \qquad
    \inferrule{\EL{\gamma, x_1 \mapsto v_1, x_2 \mapsto \sundefined}{t_1}{v}}{\ELF{\gamma}{x_1}{x_2}{t_1}{t_2}{\scons{v_1}{\sundefined}}{v}}{\rulename{L-Foldr-Undefined}} \\
    \inferrule{\ELF{\gamma}{x_1}{x_2}{t_1}{t_2}{\sthunk{v_2}}{v_3} \\
      \EL{\gamma, x_1 \mapsto v_1, x_2 \mapsto
        \sthunk{v_3}}{t_1}{v}}{\ELF{\gamma}{x_1}{x_2}{t_1}{t_2}{\scons{v_1}{v_2}}{v}}{\rulename{L-Foldr-Result}}
  \end{mathpar}
\end{center}
\caption{The semantic rules of $L$ that are related to the \TThunkZ{} datatype,
  the choice operator~(\texttt{?}), and free variables.\ly{Are there too many
    rules here?}}\label{fig:logic-semantics}
\end{figure}

\subsection{The Clairvoyance Translation}\label{subsec:clairvoyance-translation}

% \ly{This section describes properties of the clairvoyance translation, but where
%   is the clairvoyance translation?}

% We now define some constructs on the terms in \( \LL \).  For each construct \( c \), we provide typing rules and a set of evaluation rules that \emph{completely characterize} \( c \).  By ``completely characterize'', we mean the following.

% \begin{itemize}
% \item \emph{Soundness:} Each rule is valid.
% \item \emph{Completeness:} Any derivation of the form \( \EL{\gamma}{c}{v} \) must have the form of one of the elimation rules for that construct.  In particular, there must be some elimination rule whose hypotheses are all true.  This allows us to ``invert'' judgments of the form \( \EL{\gamma}{c}{v} \).
% \end{itemize}

% \subsubsection{\texttt{return}}

% \[
%   \texttt{\treturn{t}} = \tpair{t}{0}
% \]

% \begin{center}
%   \begin{prooftree}
%     \hypo{\TL{\Gamma}{t}{A}}
%     \infer1{\TL{\Gamma}{\treturn{t}}{\TM{A}}}
%   \end{prooftree}
% \end{center}

% \begin{center}
%   \begin{prooftree}
%     \hypo{\EL{\gamma}{t}{v}}
%     \infer1{\EL{\gamma}{\treturn{t}}{\spair{v}{0}}}
%   \end{prooftree}
% \end{center}

% \subsubsection{bind}

% \[
%   \texttt{do $x$ <- $t_1$; $t_2$} = \tlet{\tpair{x}{y_1}}{t_1}{\tlet{\tpair{y_2}{y_3}}{t_2}{\tpair{y_2}{y_1 + y_3}}}
% \]

% \begin{center}
%   \begin{prooftree}
%     \hypo{\TL{\Gamma}{t_1}{\TM{A}}}
%     \hypo{\TL{\Gamma, x : A}{t_2}{\TM{B}}}
%     \infer2{\TL{\Gamma}{\texttt{do $x$ <- $t_1$; $t_2$}}{\TM{B}}}
%   \end{prooftree}
% \end{center}

% \subsubsection{\texttt{transpose}}

% \[
% \begin{array}{l@{}l@{}l@{}}
%   \texttt{transpose \(t\)} = \texttt{case \(t\) of } & \texttt{thunk \(y_1\) -> do $y_2$ <- $y_1$; return (\tthunk{y_2})} \\ & \texttt{undefined -> return undefined}
% \end{array}
% \]

% \begin{center}
%   \begin{prooftree}
%     \hypo{\TL{\Gamma}{t}{\TThunk{(\TM{A})}}}
%     \infer1{\TL{\Gamma}{\ttranspose{t}}{\TM{(\TThunk{A})}}}
%   \end{prooftree}
% \end{center}

% \begin{center}
%   \begin{prooftree}
%     \hypo{\EL{\gamma}{t}{\sundefined}}
%     \infer1{\EL{\gamma}{\ttranspose{t}}{\spair{\sundefined}{0}}}
%   \end{prooftree}
% \end{center}

% \begin{center}
%   \begin{prooftree}
%     \hypo{\EL{\gamma}{t}{\sthunk{\spair{v}{c}}}}
%     \infer1{\EL{\gamma}{\ttranspose{t}}{\spair{\sthunk{v}}{c}}}
%   \end{prooftree}
% \end{center}

% \subsubsection{\texttt{foldM}}

% \[
%   \texttt{\tfoldM{x_1}{x_2}{t_1}{t_2}{t_3} = \tfoldr{x_1}{y_2}{\texttt{do \(x_2\) <- \ttranspose{y_2}; \(t_1\)}}{t_2}{t_3}} .
% \]

% \begin{center}
%   \begin{prooftree}
%     \hypo{\TL{\Gamma, x_1 : A, x_2 : \TThunk{B}}{t_1}{\TM{B}}}
%     \hypo{\TL{\Gamma}{t_2}{\TM{B}}}
%     \hypo{\TL{\Gamma}{t_3}{\TList{A}}}
%     \infer3{\TL{\Gamma}{\texttt{foldM(\(\lambda x_1 x_2\).\(t_1\), \(t_2\), \(t_3\))}}{\TM{B}}}
%   \end{prooftree}
% \end{center}

% To characterize the reduction behavior of \( \rulename{L-FoldM} \), we define yet another auxiliary relation:
% \[
%   \ELFM{\gamma}{x_1}{x_2}{t_1}{t_2}{v}{u} \equiv \ELF{\gamma}{x_1}{y_2}{\texttt{do \(x_2\) <- \ttranspose{y_2}; \(t_1\)}}{t_2}{v}{u} .
% \]

% \begin{lemma}
%   The following rules completely characterize \( \tfoldMZ \).
%   \begin{center}
%     \begin{prooftree}
%       \hypo{\ELFM{\gamma}{x_1}{x_2}{t_1}{t_2}{v}{u}}
%       \infer1[\rulename{L-FoldM}]{\EL{\gamma}{\tfoldM{x_1}{x_2}{t_1}{t_2}{v}}{u}}
%     \end{prooftree}
%   \end{center}\begin{center}
%     \begin{prooftree}
%       \hypo{\EL{\gamma}{t_2}{u}}
%       \infer1[\rulename{L-FoldM-Nil}]{\ELFM{\gamma}{x_1}{x_2}{t_1}{t_2}{\varepsilon}{u}}
%     \end{prooftree}
%   \end{center}
%   \begin{center}
%     \begin{prooftree}
%       \hypo{\EL{\gamma, x_1 \mapsto v_1, x_2 \mapsto \sundefined}{t_1}{u}}
%       \infer1[\rulename{L-FoldM-Undefined}]{\ELFM{\gamma}{x_1}{x_2}{t_1}{t_2}{\scons{v_1}{\sundefined}}{u}}
%     \end{prooftree}
%   \end{center}
%   \begin{center}
%     \begin{prooftree}
%       \hypo{\ELFM{\gamma}{x_1}{x_2}{t_1}{t_2}{v_2}{\spair{v_3}{c_1}}}
%       \hypo{\EL{\gamma, x_1 \mapsto v_1, x_2 \mapsto \sthunk{v_2}}{t_1}{\spair{v_4}{c_2}}}
%       \infer2[\rulename{L-FoldM-Thunk}]{\ELFM{\gamma}{x_1}{x_2}{t_1}{t_2}{\scons{v_1}{\sthunk{v_2}}}{\spair{v_4}{c_1 + c_2}}}
%     \end{prooftree}
%   \end{center}
% \end{lemma}
% \ifextended
% \begin{proof}
%   We first prove validity.
%   \newpage
%   \begin{landscape}
%     \begin{itemize}
%     \item \textbf{\rulename{L-FoldM}.}
%       \begin{center}
%         \begin{prooftree}
%           \hypo{\ELF{\gamma}{x_1}{y_2}{\texttt{do \(x_2\) <- \ttranspose{y_2}; \(t_1\)}}{t_2}{v}{u}}
%           \infer1[\rulename{L-Foldr}]{\EL{\gamma}{\tfoldr{x_1}{y_2}{\texttt{do \(x_2\) <- \ttranspose{y_2}; \(t_1\)}}{t_2}{v}}{u}}
%         \end{prooftree}
%       \end{center}
%     \item \textbf{\rulename{L-FoldM-Nil}.}
%       \begin{center}
%         \begin{prooftree}
%           \hypo{\EL{\gamma}{t_2}{u}}
%           \infer1[\rulename{L-Foldr-Nil}]{\ELF{\gamma}{x_1}{y_2}{\texttt{do \(x_2\) <- \ttranspose{y_2}; \(t_1\)}}{t_2}{\varepsilon}{u}}
%         \end{prooftree}
%       \end{center}
%     \item \textbf{\rulename{L-FoldM-Undefined}.}
%       \begin{center}
%         \begin{prooftree}
%           \infer0[\rulename{L-Var}]{\EL{\gamma, x_1 \mapsto v_1, y_2 \mapsto \bot}{y_2}{\bot}}
%           \infer1[\rulename{L-Transpose-Undefined}]{\EL{\gamma, x_1 \mapsto v_1, y_2 \mapsto \bot}{\ttranspose{y_2}}{\spair{\bot}{0}}}
%           \hypo{\EL{\gamma, x_1 \mapsto v_1, y_2 \mapsto \bot, x_2 \mapsto \sundefined}{t_1}{u}}
%           \infer2[\rulename{L-Bind}]{\EL{\gamma, x_1 \mapsto v_1, y_2 \mapsto \bot}{\texttt{do \(x_2\) <- \ttranspose{y_2}; \(t_1\)}}{u}}
%           \infer1[\rulename{L-Foldr-Undefined}]{\ELF{\gamma}{x_1}{y_2}{\texttt{do \(x_2\) <- \ttranspose{y_2}; \(t_1\)}}{t_2}{\scons{v_1}{\sundefined}}{u}}
%         \end{prooftree}
%       \end{center}
%     \item \textbf{\rulename{L-FoldM-Thunk}.}  By induction, \( \ELF{\gamma}{x_1}{y_2}{\texttt{do \(x_2\) <- \ttranspose{y_2}; \(t_1\)}}{t_2}{v_2}{v_3} \).  Then,
%       \begin{center}
%         \begin{prooftree}
%           \hypo{\ELF{\gamma}{x_1}{y_2}{\texttt{do \(x_2\) <- \ttranspose{y_2}; \(t_1\)}}{t_2}{v_2}{v_3}}
%           \infer0[\rulename{L-Var}]{\EL{\gamma, x_1 \mapsto v_1, y_2 \mapsto \sthunk{\spair{v_3}{c_1}}}{y_2}{\sthunk{\spair{v_3}{c_1}}}}
%           \infer1[\rulename{L-Transpose-Thunk}]{\EL{\gamma, x_1 \mapsto v_1, y_2 \mapsto \sthunk{\spair{v_3}{c_1}}}{\ttranspose{y_2}}{\spair{\sthunk{v_3}}{c_1}}}
%           \hypo{\EL{\gamma, x_1 \mapsto v_1, y_2 \mapsto \sthunk{\spair{v_3}{c_2}}, x_2 \mapsto \sthunk{v_2}}{t_1}{\spair{v_4}{c_2}}}
%           \infer2[\rulename{L-Bind}]{\EL{\gamma, x_1 \mapsto v_1, y_2 \mapsto \sthunk{\spair{v_3}{c_2}}}{\texttt{do \(x_2\) <- \ttranspose{y_2}; \(t_1\)}}{\spair{v_4}{c_1 + c_2}}}
%           \infer2[\rulename{L-Foldr-Thunk}]{\ELF{\gamma}{x_1}{y_2}{\texttt{do \(x_2\) <- \ttranspose{y_2}; \(t_1\)}}{t_2}{\scons{v_1}{\sthunk{v_2}}}{\spair{v_4}{c_1 + c_2}}}
%         \end{prooftree}
%       \end{center}
%     \end{itemize}
%     To prove completeness, repeatedly apply inversion principles to find that any derivation of \( \tfoldM{x_1}{x_2}{t_1}{t_2}{v}{u} \) must have the form of one of the preceding derivations.
%   \end{landscape}
% \end{proof}
% \fi

\begin{figure}[t]
  \[
    \begin{array}{rcl}
      \tassert{t_1}{t_2} &= &\tif{t_1}{t_2}{\tfail} \\
      \treturn{t} &= &\tpair{t}{0} \\
      \tdo{x}{t_1}{t_2} &= &\tlet{y_1}{t_1}{\relax} \\ &&\tlet{x}{\tfst{t_1}}{\relax} \\ &&\tlet{y_2}{t_2}{\relax} \\ && \tpair{\tfst{y_2}}{\tadd{\tsnd{y_1}}{\tsnd{y_2}}} \\
      \tfoldM{x_1}{x_2}{t_1}{t_2}{t_3} &= &\tfoldr{x_1}{y_2}{\ttranspose{y_2}}{t_2}{t_3} \\
      \ttranspose{t} &= &\texttt{case \(t\) of } \\ &&\texttt{thunk \(y_1\) -> (\tdo{y_2}{y_1}{\treturn{(\tthunk{y_2})}});} \\ &&\texttt{undefined -> return undefined}
    \end{array}
  \]
\caption{Derived term constructs in \(\LL\).}
\end{figure}

\begin{figure}[t]
  \centering
  \begin{mathpar}
    \inferrule{\TL{\Gamma}{t_1}{\TBool} \\ \TL{\Gamma}{t_2}{A}}{\TL{\Gamma}{\tassert{t_1}{t_2}}{A}} \\
    \inferrule{\TL{\Gamma}{t}{A}}{\TL{\Gamma}{\treturn{t}}{\TM{A}}} \qquad
    \inferrule{\TL{\Gamma}{t_1}{\TM{A}} \\ \TL{\Gamma, x : A}{t_2}{\TM{B}}}{\TL{\Gamma}{\texttt{\tdo{x}{t_1}{t_2}}}{\TM{B}}} \\
    \inferrule{\TL{\Gamma, x_1 : A, x_2 : \TThunk{B}}{t_1}{\TM{B}} \\ \TL{\Gamma}{t_2}{\TM{B}} \TL{\Gamma}{t_3}{\TList{A}}}{\TL{\Gamma}{\tfoldM{x_1}{x_2}{t_1}{t_2}{t_3}}{\TM{B}}}
  \end{mathpar}
  \caption{Typing rules for the derived constructs.}
\end{figure}

\begin{figure}[t]
  \begin{align*}
    \transC{x} &= \treturn{x} \\
    \transC{\tlet{x}{t_1}{t_2}} &= \tdo{x}{\transC{t_1}}{\transC{t_2}} \\
    \transC{\ttrue{}} &= \treturn{\ttrue{}} \\
    \transC{\tfalse{}} &= \treturn{\tfalse{}} \\
    \transC{\tif{t_1}{t_2}{t_3}} &= \tdo{y}{t_1}{\tif{y}{\transC{t_2}}{\transC{t_3}}} \\
    \transC{\tnil{}} &= \treturn{\tnil{}} \\
    \transC{\tcons{t_1}{t_2}} &= \tdo{y_1}{\transC{t_1}}{\tdo{y_2}{\transC{t_2}}{\treturn{(\tcons{y_1}{y_2})}}} \\
    \transC{\tfoldr{x_1}{x_2}{t_1}{t_2}{t_3}} &= \tdo{y}{\transC{t_3}}{\tfoldM{x_1}{x_2}{\transC{t_1}}{\transC{t_2}}{y}} \\
    \transC{\tlazy{t}} &= \tchoose{\treturn{\tundefined}}{\tdo{y}{t_1}{\treturn(\tthunk{y})}} \\
    \transC{\tforce{t}} &= \tdo{y_1}{t}{\tcase{y_1}{y_2}{\treturn{y_2}}{\tfail}}
  \end{align*}
  \caption{The clairvoyance translation \(\transC{-}\) from \(\LB\) to \(\LL\).}
  \label{fig:transC}
\end{figure}

We define a translation \(\transC{-}\) from \(\LB\) to \(\LL\) that is intended to ``behave like'' clairvoyance semantics (\cref{fig:transC}).

\ly{In the lemma below, should the two $\Gamma$s be the identical?}\nc{Yes: as
  we've defined the semantics, the set of types in \(\LL\) is a superset of the
  set of types in \(\LB\), so any typing context for \(\LB\) is also a valid
  typing context for \(\LL\). Likewise, the set of values in \(\LL\) is a
  superset of the set of partial values in \(\LB\), so any evaluation context
  that is valid for clairvoyance semantics is also a valid evaluation context
  for \(\LL\).}

\begin{lemma}
  If \( \TB{\Gamma}{t}{A} \), then \( \TL{\Gamma}{\transC{t}}{\TM{A}} \).
\end{lemma}

\subsection{The Demand Translation}\label{subsec:demand-translation}

We now define another translation from \( \LB \) to \( \LL \) that is intended
to simulate demand semantics. Given \( \TB{\Gamma}{t}{A} \), the demand
translation of \( t \) introcuces new logic variables with types corresponding
to the types in \( \Gamma \); consequently, \( \Gamma \) must become a parameter
to the translation. We therefore denote the demand translation by
\( \transD{t}{\Gamma} \).

\begin{figure}[t]
  \[
    \begin{array}{rl}
      \transD{t}{x_1 : A_1, \ldots, x_n : A_n} = &\tlet{x'_1}{\tfree{A_1}}{\relax} \\
                                                 &\tassert{\tle{x'_1}{x_1}}{\relax} \\
                                                 &\multicolumn{1}{c}{\vdots} \\
                                                 &\tlet{x'_n}{\tfree{A_n}}{\relax} \\
                                                 &\tassert{\tle{x'_n}{x_n}}{\relax} \\
                                                 &\tlet{u}{\transC{t'}}{\relax} \\
                                                 &\tassert{\teq{x}{\tfst{u}}}{\relax} \\
                                                 &\tpair{\texttt{(\( x'_1 \),\ \( \ldots \),\ \( x'_n \))}}{\tsnd{u}} \\
      \multicolumn{2}{c}{\text{where \( t' = t[x'_1/x_1, \ldots, x'_n/x_n] \).}}
    \end{array}
  \]
  \caption{The demand translation.}
\end{figure}

Since the demand translation is intended to simulate demand semantics, it must
evaluate to ``the same sort of thing'' as demand semantics, \ie an environment
of approximations. To that end, we define an operation \( \verts{-} \) that
turns contexts into types:
\[
  \verts{x_1 : A_1, \ldots, x_n : A_1} = \TProd{A_1}{\TProd{\cdots}{A_n}} ,
\]
with the degenerate case \( \verts{\varepsilon} = \TUnit \).

\begin{lemma}
  If \( \TB{\Gamma}{t}{A} \), then \( \TL{\Gamma, x : A}{\transD{t}{\Gamma}}{\TM{\verts{\Gamma}}} \).
\end{lemma}

\subsection{Example: Revisiting \texttt{take}}

\ly{This section talks about how one can apply the translation to the example of
  \texttt{take} and \texttt{append}.}
